\documentclass[11pt,letterpaper]{article}
\usepackage[T1]{fontenc}
\usepackage{lmodern}
\usepackage[letterpaper,left=0.65in,right=0.65in,top=0.55in,bottom=0.65in,footskip=0.31in]{geometry}
\usepackage{amsmath}
\usepackage{tikz}
\usetikzlibrary{arrows.meta}
\usepackage{fancyhdr}

\setlength{\parindent}{0pt}
\setlength{\parskip}{0pt}
\pagestyle{fancy}
\fancyhf{}
\fancyfoot[L]{\footnotesize EE 102}
\fancyfoot[R]{\footnotesize\thepage}
\renewcommand{\headrulewidth}{0pt}
\renewcommand{\footrulewidth}{0.3pt}
\newcommand{\problem}[2]{\par\noindent{\bfseries#1.\enspace #2}\par\smallskip}
\newcommand{\blank}[1]{\rule{#1}{0.4pt}}

% The same axes and dimensions are used in all nine boxes.
% A panel's third argument contains its trace. An empty argument leaves blank axes.
\tikzset{
  trace/.style={draw=black,line width=1.15pt,line cap=round,line join=round},
  axis/.style={draw=black,line width=0.45pt,-{Stealth[length=3.7pt]}},
  ticklabel/.style={font=\scriptsize,inner sep=0.7pt}
}
\newcommand{\panel}[3]{%
  \begin{tikzpicture}[x=0.51cm,y=1.00cm,baseline=(current bounding box.north)]
    \path[use as bounding box] (-5.9,-0.58) rectangle (5.6,3.52);
    \draw[black] (-5.9,-0.58) rectangle (5.6,3.52);
    \node[anchor=west,inner sep=0pt,font=\normalsize] at (-5.45,3.08) {$#1$};
    \node[anchor=east,inner sep=0pt,font=\scriptsize] at (5.15,3.08) {#2};
    \draw[axis] (-4.85,0) -- (5.08,0);
    \draw[axis] (-4.85,0) -- (-4.85,2.57);
    \foreach \w in {-4,-2,0,2,4}{
      \draw[black,line width=0.4pt] (\w,-0.05) -- (\w,0.05);
      \node[ticklabel,anchor=north] at (\w,-0.075) {$\w$};
    }
    \foreach \a in {1,2}{
      \draw[black,line width=0.4pt] (-4.94,\a) -- (-4.76,\a);
      \node[ticklabel,anchor=east] at (-4.99,\a) {$\a$};
    }
    \node[ticklabel,anchor=north] at (5.15,-0.075) {$\omega$};
    #3
  \end{tikzpicture}%
}
\newcommand{\sketchrow}[3]{\par\smallskip\noindent\hbox to \linewidth{#1\hfil#2\hfil#3}\par}
\newcommand{\workbox}[1]{\par\noindent\fbox{\parbox[c][#1][c]{\dimexpr\linewidth-2\fboxsep-2\fboxrule\relax}{\strut}}\par}

\begin{document}
{EE 102}\hfill Name \blank{2.45in}\par
\vspace{0.09in}
Audio Processing Using Fourier Transforms\par
\vspace{0.10in}
Let $t$ denote time in seconds and $\omega$ angular frequency in rad/s.
The continuous-time Fourier transform and its inverse are
\[
  X(\omega)=\int_{-\infty}^{\infty}x(t)e^{-j\omega t}\,dt,
  \qquad
  x(t)=\frac{1}{2\pi}\int_{-\infty}^{\infty}X(\omega)e^{j\omega t}\,d\omega .
\]
For an LTI system, let $x(t)$ be the input, $h(t)$ its impulse response, and $y(t)$ the output.
Write $H(\omega)$ and $Y(\omega)$ for the transforms of $h(t)$ and $y(t)$. Then
\[
  y(t)=(h*x)(t)=\int_{-\infty}^{\infty}x(\tau)h(t-\tau)\,d\tau
  \quad\overset{\mathcal F}{\longleftrightarrow}\quad
  Y(\omega)=H(\omega)X(\omega) .
\]

\problem{1}{Sketch the filter}
Sketch $H(\omega)$ for the input and output shown.
\sketchrow
  {\panel{X(\omega)}{Given}{\draw[trace] (-4.75,0)--(-4,0)--(0,2)--(4,0)--(4.75,0);}}
  {\panel{H(\omega)}{Your sketch}{}}
  {\panel{Y(\omega)}{Given}{\draw[trace] (-4.75,0)--(-2,0)--(-2,1)--(0,2)--(2,1)--(2,0)--(4.75,0);}}

\vfill
\problem{2}{Sketch the output}
Sketch $Y(\omega)$ for the input and filter shown.
\sketchrow
  {\panel{X(\omega)}{Given}{\draw[trace] (-4.75,0)--(-4,0)--(-2,2)--(2,2)--(4,0)--(4.75,0);}}
  {\panel{H(\omega)}{Given}{\draw[trace] (-4.75,1)--(-2,1)--(-2,0)--(2,0)--(2,1)--(4.75,1);}}
  {\panel{Y(\omega)}{Your sketch}{}}

\vfill
\problem{3}{Recover the input}
Sketch $X(\omega)$ for the filter and output shown.
\sketchrow
  {\panel{X(\omega)}{Your sketch}{}}
  {\panel{H(\omega)}{Given}{\draw[trace] (-4.75,1)--(-2,1)--(-2,2)--(2,2)--(2,1)--(4.75,1);}}
  {\panel{Y(\omega)}{Given}{\draw[trace] (-4.75,0)--(-4,0)--(-4,2)--(4,2)--(4,0)--(4.75,0);}}

\vspace*{0.08in}
\newpage
{\bfseries EE 102}\hfill{Audio Processing Using Fourier Transforms}\par
\vspace{0.20in}
\problem{4}{Find the output in the frequency domain}
Let $a>0$. The frequency-domain representations of the input and filter are
\[
  X(\omega)=\frac{1}{a+j\omega},
  \qquad
  H(\omega)=
  \begin{cases}
    1, & |\omega|<10\ \mathrm{rad/s},\\[2pt]
    0, & |\omega|\ge 10\ \mathrm{rad/s}.
  \end{cases}
\]
Write $Y(\omega)$ as a piecewise function for all $\omega$.\par\smallskip
\workbox{1.30in}

\vspace{0.28in}
\textbf{Hands-on Activity: Guitar Fourier Lab}\par\smallskip
Use the guitar recording in the app. The filter editor uses frequency $f$ in Hz, with $\omega=2\pi f$.\par
\vspace{0.16in}
\problem{5}{Separate the low and high parts of the clipped guitar}
Choose a clipping threshold and a cutoff frequency. Design one filter to hear the sound below the cutoff and another to hear the sound above it.\par
\vspace{0.16in}
Clipping threshold \blank{0.65in}\,\% \hfill Cutoff frequency \blank{1.0in}\,Hz\hspace{0.25in}\par
\vspace{0.18in}
What do you hear below the cutoff?\par\vspace{0.17in}\blank{\linewidth}\par
\vspace{0.16in}
What do you hear above the cutoff?\par\vspace{0.17in}\blank{\linewidth}\par

\vspace{0.26in}
\problem{6}{Single out one note}
Set the clipping threshold to $100\%$ and isolate one note (your favorite note in this guitar song file) with a narrow filter (design that filter!).\par
\vspace{0.18in}
Fundamental frequency \blank{0.65in}\,Hz \hfill
Passband \blank{0.55in} to \blank{0.55in}\,Hz\hspace{0.25in}\par
\vspace{0.17in}
When do you hear the note in the filtered recording?\par\vspace{0.18in}\blank{\linewidth}\par

\end{document}
